\section{Discussion: when does a reusable smoothness interface pay off?}\label{sec:discussion}
The main distinction is between \emph{a mathematical coordinate} and \emph{a computational policy}. The normalized smoothness map $S_d$ is an exact reparameterization for any $1\le d<n$ with fixed numeric $\mu$. It is valuable to users because the range and direction of smoothness remain interpretable across $d$ even though penalties differ radically in magnitude. Nothing in this argument implies that optimizing in $s$ must take less time than in $\lambda$; in fact, evaluating $s\mapsto\lambda$ at every optimizer call introduces extra computation. The measured advantage appears when the same candidate grid and spectral information are reused across calls.

Figure~\ref{fig:coordinates} and Table~\ref{tab:criteria} use $d=3$ and demonstrate that native and premapped full-$S$ searches have identical statistical candidates. An independently sampled log-$\lambda$ grid was included in the original experiment solely as a conventional computational control; because it selects among different interior smoothness levels, we remove it from the main display rather than treating it as a competing interpretable interface. Bounded optimizers likewise take different search paths and do not certify global optimality.

The best numerical backend depends on the \emph{actual requested output}. A banded factorization can be very economical for one fitted curve, whereas amortizing one dense SVD can accelerate repeated spectral GCV/ML scores. Exact CV leverages and specialized trace computation can reverse rankings yet again \citep{weinert2007,biessy2026}. High-$s$, large-$d$ conditioning matters too. The independent accuracy companion found that full SVD with shift-free reconstruction is especially reliable in the tested extreme regime. However, that finding should not be back-projected as a universal wall-time speed ranking, nor should older valid low-$s$ benchmark cells be discarded.

\subsection{What the historical factorial does not establish}
Although the archived Windows source is a complete factorial \emph{within its declared axes}, it covers $d=0,\ldots,5$ rather than $1,\ldots,6$ and only fixed $\mu=0$ versus fixed $\mu=0.15$. It does not test $\widehat\mu=\operatorname{mean}(D_dy)$ as a separately timed policy. Nor does a two-seed, $n=24,48$ matrix establish asymptotic constants or algorithmic scaling laws at arbitrary $n$. We therefore report the big-$O$ analysis independently of the machine measurements. Future work should add $d=6$, data-derived and jointly estimated reference means, wider $n$, additional independent timing sessions and statistical evaluation of complete-workflow validity.

For a truly complete cross-backend selection comparison, the selected-inverse/trace approach of \citet{weinert2007} is essential prior art. Additional workloads should distinguish an observed-difference plug-in from fixed $\mu$, conditional LOOCV from chronological forecast validation, independent-$\lambda$ from matched-$s$ candidate grids, and numerical failures from statistically undefined criterion values. The residual-based BIC interpolation limit should be restricted explicitly if included as a selector, as detailed by the companion reliability study.

\subsection{Toward one combined manuscript}
We propose ultimately combining: (i) the mathematical $s\leftrightarrow\lambda$ bijection and fixed/nonzero-$\mu$ formulas; (ii) the present algorithmic cost models, matched-coordinate parity and amortized runtime evidence; and (iii) independent high-precision solver validation and finite-sample latent-trend RMSE from the numerical companion. The comparison should use a \emph{joint accuracy--time admissible frontier}: a solver must deliver valid equivalent criterion values over the \emph{complete} candidate set before its timing is ranked. This avoids favoring fast, inaccurate or incomplete high-order factorizations and clarifies whether a cached spectral default is the best robust option for a specified workload.

\section{Conclusions}
A normalized smoothness index is a mathematically interpretable interface for quadratic finite-difference PLS at any $1\le d<n$, independent of the fixed mean of differences. The full singular spectrum makes the map invertible, permits reusable candidate grids and enables amortized criterion scoring. The revised Windows analysis compares an identical complete $S$ grid at $d=3$, measuring 10.050\,ms for repeated native inversion versus 3.767\,ms with one cached grid across five criteria (2.668-fold). The historical $d=2$ case separately gives a 2.823-fold reduction. Neither ratio is a universal speedup, and only two backends completed the difficult $d=3$ full-grid GCV tasks validly. Other complete-fit benchmarks show why banded solvers remain valuable, and no universal fastest backend follows. This computational companion establishes the complexity and runtime half of a planned combined accuracy--efficiency study, with the remaining design gaps explicitly identified.

\section*{Data and code availability}
Immutable original timings, seed-level results, validations, scripts and environment metadata reside under \url{https://github.com/heri-espino/trendestimation/tree/main/paper/benchmarks/}. The original source paths and their restrictions are catalogued in this manuscript's \texttt{SOURCE\_AUDIT.md}. This paper and its TeX-native figures reside in \texttt{paper/csda-computational-costs/}. No archived numerical results are overwritten or reclassified.

\section*{Declarations}
Provisional authorship and affiliation are shown for drafting only. Author order, funding, conflicts and contributor statements require author approval before submission.
